\section*{Capítulo \ref{ch:g11:organic-skeletons} --- Las moléculas orgánicas: esqueletos y nombres}

\begin{solution}{exo:g11:organic-skeletons:1}
Las dos tienen seis \omterm{def:g7:atoms-and-molecules:atom}{átomos} de carbono (extremos y vértices) y son
\ce{C6H14}: el hexano y el 2-metilpentano son \omterm{def:g11:organic-skeletons:isomer}{isómeros}.
\end{solution}

\begin{solution}{exo:g11:organic-skeletons:2}
Propano; hexano; 2-metilbutano.
\end{solution}

\begin{solution}{exo:g11:organic-skeletons:3}
\ce{CH3-CH2-CH3}; \chemfig{CH_3-CH(-[2]CH_3)-CH_3};
\chemfig{H-C(-[2]H)(-[6]H)-C(-[2]H)(-[6]H)-H}.
\end{solution}

\begin{solution}{exo:g11:organic-skeletons:4}
La primera, el 3-metilpentano: 6 \omterm{def:g7:atoms-and-molecules:atom}{átomos} de carbono; sus extremos
\ce{CH3} (tres) llevan 9 \omterm{def:g7:atoms-and-molecules:atom}{átomos} de hidrógeno, sus dos vértices \ce{CH2},
4, y su punto de ramificación \ce{CH}, 1: \ce{C6H14}. El ciclohexano: 6
vértices, cada uno \ce{CH2}: \ce{C6H12}.
\end{solution}

\begin{solution}{exo:g11:organic-skeletons:5}
\ce{C8H18}. $2n + 2 = 20$ da $n = 9$: el nonano, \ce{C9H20}.
\end{solution}

\begin{solution}{exo:g11:organic-skeletons:6}
\setchemfig{atom sep=1.6em}
3-metilhexano \chemfig{-[:30]-[:-30]-[:30](-[:90])-[:-30]-[:30]-[:-30]},
\ce{C7H16}; 2,2-dimetilbutano \chemfig{-[:0](-[:90])(-[:-90])-[:0]-[:30]},
\ce{C6H14}; 3-etilpentano
\chemfig{-[:30]-[:-30](-[:-90]-[:-30])-[:30]-[:-30]}, \ce{C7H16}.
\end{solution}

\begin{solution}{exo:g11:organic-skeletons:7}
El pentano, el 2-metilbutano y el 2,2-dimetilpropano, dibujados en la
figura del capítulo.
\end{solution}

\begin{solution}{exo:g11:organic-skeletons:8}
El ``2-etilbutano'' tiene una cadena más larga de cinco carbonos que
pasa por el grupo etilo: 3-metilpentano. El ``4-metilpentano'' está
numerado desde el extremo equivocado: 2-metilpentano. El
``1-metilbutano'' es simplemente una cadena de cinco carbonos: pentano.
\end{solution}

\begin{solution}{exo:g11:organic-skeletons:9}
2,2-dimetilpropano (\qty{9.5}{\celsius}) < 2-metilbutano
(\qty{27.8}{\celsius}) < pentano (\qty{36.1}{\celsius}). Los mismos
\omterm{def:g7:atoms-and-molecules:atom}{átomos} están dispuestos de forma cada vez más compacta: menos contacto
entre \omterm{def:g7:atoms-and-molecules:molecule}{moléculas}, \omterm{def:g11:polarity-and-cohesion:van-der-waals}{interacciones de van der Waals} más débiles y
temperatura de ebullición más baja.
\end{solution}

\begin{solution}{exo:g11:organic-skeletons:10}
Al cerrar la cadena en un anillo se forma un enlace \ce{C-C} más, que
ocupa el lugar de dos enlaces \ce{C-H} (uno en cada extremo). \omterm{def:g11:organic-skeletons:formulas}{Fórmulas
moleculares} distintas: no son \omterm{def:g11:organic-skeletons:isomer}{isómeros}. El ciclohexano no es un \omterm{def:g11:organic-skeletons:alkane}{alcano}
(su esqueleto es un anillo; su fórmula no es \ce{C_{n}H_{2n+2}}).
\end{solution}

\begin{solution}{exo:g11:organic-skeletons:11}
El butano y el 2-metilpropano (\ce{C4H10}). El propano es \ce{C3H8}, y
el ciclobutano, \ce{C4H8}.
\end{solution}

\begin{solution}{exo:g11:organic-skeletons:12}
\setchemfig{atom sep=1.5em}
\begin{center}
\small
\begin{tabular}{ccc}
\chemfig{-[:30]-[:-30]-[:30]-[:-30]-[:30]} &
\chemfig{-[:30](-[:90])-[:-30]-[:30]-[:-30]} &
\chemfig{-[:30]-[:-30](-[:-90])-[:30]-[:-30]} \\
hexano & 2-metilpentano & 3-metilpentano \\[6pt]
\chemfig{-[:0](-[:90])(-[:-90])-[:0]-[:30]} &
\chemfig{-[:30](-[:90])-[:-30](-[:-90])-[:30]} & \\
2,2-dimetilbutano & 2,3-dimetilbutano &
\end{tabular}
\end{center}
\end{solution}

\begin{solution}{exo:g11:organic-skeletons:13}
Dos zigzags largos pueden ponerse uno al lado del otro a lo largo de
toda su longitud, con muchos \omterm{def:g7:atoms-and-molecules:atom}{átomos} en contacto; dos bolas solo se
tocan en una superficie pequeña. Las \omterm{def:g11:polarity-and-cohesion:van-der-waals}{interacciones de van der Waals}
actúan entre \omterm{def:g7:atoms-and-molecules:atom}{átomos} en contacto: son mayores entre \omterm{def:g7:atoms-and-molecules:molecule}{moléculas} lineales,
que necesitan una temperatura más alta para separarse.
\end{solution}

\begin{solution}{exo:g11:organic-skeletons:14}
La cadena más larga tiene seis carbonos; numerada desde la izquierda,
los sustituyentes están en 2, 4, 4, y desde la derecha, en 3, 3, 5; la
primera diferencia (2 frente a 3) decide: 2,4,4-trimetilhexano.
\end{solution}

\begin{solution}{exo:g11:organic-skeletons:15}
\setchemfig{atom sep=1.6em}
\chemfig{-[:0](-[:90])(-[:-90])-[:-30]-[:30](-[:90])-[:-30]}: una cadena
de cinco carbonos con dos grupos metilo en el carbono 2 y uno en el
carbono 4, ocho \omterm{def:g7:atoms-and-molecules:atom}{átomos} de carbono en total; $2 \times 8 + 2 = 18$ \omterm{def:g7:atoms-and-molecules:atom}{átomos}
de hidrógeno, \ce{C8H18}, la fórmula del octano: son \omterm{def:g11:organic-skeletons:isomer}{isómeros}. \omterm{def:g11:organic-skeletons:formulas}{Fórmula
semidesarrollada}:
\chemfig{CH_3-C(-[2]CH_3)(-[6]CH_3)-CH_2-CH(-[2]CH_3)-CH_3}.
\end{solution}

\begin{solution}{pb:g11:organic-skeletons:1}
\setchemfig{atom sep=1.5em}
\textbf{1.} \ce{C4H10}: con $n = 4$, $2n + 2 = 10$.

\textbf{2.} \chemfig{H-C(-[2]H)(-[6]H)-C(-[2]H)(-[6]H)-C(-[2]H)(-[6]H)-C(-[2]H)(-[6]H)-H}.

\textbf{3.} \ce{CH3-CH2-CH2-CH3}; \chemfig{-[:30]-[:-30]-[:30]}.

\textbf{4.} \chemfig{H-C(-[2]H)(-[6]H)-C(-[6]H)(-[2,1.5]C(-[1]H)(-[3]H)(-[2]H))-C(-[2]H)(-[6]H)-H};
\chemfig{CH_3-CH(-[2]CH_3)-CH_3}; \chemfig{-[:30](-[:90])-[:-30]}.
Tiene los mismos \omterm{def:g7:atoms-and-molecules:atom}{átomos}, cuatro carbonos y diez hidrógenos, unidos en un
orden distinto.

\textbf{5.} Sí: \omterm{def:g11:organic-skeletons:isomer}{isómeros constitucionales} (esqueletos distintos).

\textbf{6.} Los tres: sus temperaturas de ebullición están por debajo
de \qty{20}{\celsius}.

\textbf{7.} A \qty{-5}{\celsius}, por debajo de su temperatura de
ebullición, el butano sigue líquido a presión normal y desprende poco
gas.

\textbf{8.} El propano (\qty{-42}{\celsius}) y el 2-metilpropano
(\qty{-11.7}{\celsius}) hierven por debajo de las temperaturas del
invierno: siguen dando mucho gas con frío.

\textbf{9.} El 2-metilpropano es ramificado, más compacto:
\omterm{def:g11:polarity-and-cohesion:van-der-waals}{interacciones de van der Waals} más débiles.

\textbf{10.} Con tres \omterm{def:g7:atoms-and-molecules:atom}{átomos} de carbono, el único esqueleto posible es la
cadena \ce{C-C-C}: cualquier carbono que se cambie de sitio da la misma
cadena.

\textbf{11.} 2-metilbutano.

\textbf{12.} La cadena se debe numerar desde el extremo más cercano al
sustituyente: 2, no 3.

\textbf{13.} 2, 3 y 5.

\textbf{14.} Heptano, \chemfig{-[:30]-[:-30]-[:30]-[:-30]-[:30]-[:-30]}.

\textbf{15.} 2-metilhexano \chemfig{-[:30](-[:90])-[:-30]-[:30]-[:-30]-[:30]}
y 3-metilhexano \chemfig{-[:30]-[:-30](-[:-90])-[:30]-[:-30]-[:30]}. El
``4-metilhexano'' es el 3-metilhexano numerado desde el extremo
equivocado.

\textbf{16.} 2,2-dimetilpentano, 2,3-dimetilpentano,
2,4-dimetilpentano y 3,3-dimetilpentano.

\textbf{17.} Sí, el 3-etilpentano. En el carbono 2, el grupo etilo daría
una cadena de seis carbonos que pasaría por él: la \omterm{def:g7:atoms-and-molecules:molecule}{molécula} sería el
3-metilhexano.

\textbf{18.} 2,2,3-trimetilbutano,
\chemfig{-[:0](-[:90])(-[:-90])-[:0](-[:90])-[:0]}.

\textbf{19.} $1 + 2 + 4 + 1 + 1 = 9$.

\textbf{20.} El heptano tiene 9 \omterm{def:g11:organic-skeletons:isomer}{isómeros constitucionales}, incluido él
mismo.
\end{solution}
